\documentclass{article}

\usepackage[utf8]{inputenc}     % pdfLaTeX
\usepackage[T1]{fontenc}
\usepackage[magyar]{babel}
\usepackage{amsmath}

% Set page size and margins
% Replace `letterpaper' with `a4paper' for UK/EU standard size
\usepackage[a4paper,top=2cm,bottom=2cm,left=3cm,right=3cm,marginparwidth=1.75cm]{geometry}
%for arrows
\usepackage{textcomp}

\title{Kötőszók}
\author{Lengyel Zoltán Péter\thanks{Csukás Ibolya tanítása alapján}}
\date{Legutóbb frissítve: 2026. Március 20.}

\begin{document}

\maketitle
\tableofcontents
\newpage

\begin{center}
    \section{Kötőszók}
\end{center}

\begin{itemize}
    \item Mondaton belüli nyelvtani és logikai viszonyokat jelölnek
    \item Összekapcsolhatnak szavakat, de akár a szavaknál nagyobb nyelvi egységeket is
\end{itemize}

3 fajtájuk van:
\begin{itemize}

    \item Egyes kötőszavak: 
    \begin{itemize}
        \item ha, hogy, mert, hanem, ezért, azaz
    \end{itemize}
    
    \item Páros kötőszavak 
    \begin{itemize}
        \item is-is (szép is, jó is)
        \item sem-sem (sem apró, sem nagy)
        \item vagy-vagy (vagy kék, vagy piros)
    \end{itemize}
    
    \item Hármas kötőszó 
    \begin{itemize}
        \item nemcsak-hanem-is (nemcsak szép, hanem okos is)
    \end{itemize}
    
\end{itemize}

\end{document}